
\documentclass[11pt]{article}

\usepackage[margin=1in]{geometry}
\usepackage{amsmath, amssymb, amsthm, mathtools}
\usepackage{booktabs}
\usepackage{enumitem}
\usepackage{hyperref}
\usepackage{natbib}
\usepackage{graphicx}

\hypersetup{
  colorlinks=true,
  linkcolor=blue,
  citecolor=blue,
  urlcolor=blue
}

\newtheorem{assumption}{Assumption}
\newtheorem{proposition}{Proposition}
\newtheorem{theorem}{Theorem}
\newtheorem{corollary}{Corollary}
\newtheorem{lemma}{Lemma}
\theoremstyle{remark}
\newtheorem{remark}{Remark}

\newcommand{\E}{\mathbb{E}}
\newcommand{\Var}{\mathrm{Var}}
\newcommand{\Cov}{\mathrm{Cov}}
\newcommand{\R}{\mathbb{R}}
\newcommand{\TV}{\|\cdot\|_{\mathrm{TV}}}
\newcommand{\Wone}{\mathcal{W}_1}
\newcommand{\norm}[1]{\left\|#1\right\|}
\newcommand{\ip}[2]{\left\langle #1,#2\right\rangle}
\newcommand{\dd}{\,\mathrm{d}}

\title{Pathwise Perturbation Particle Filtering\\
with Anticipated-Shock Quadrature and Piecewise Localization}
\author{Ed \and \texttt{(coauthor to be added)}}
\date{\today}

\begin{document}
\maketitle

\begin{abstract}
We formalize a likelihood-evaluation strategy for nonlinear DSGE models based on \emph{pathwise} (a.k.a.\ extended-path) perfect-foresight solves and first-order linearizations around those paths.
At each filtering step we augment the model with an auxiliary ``anticipated-shock'' object---a finite horizon of future structural innovations---that is used to approximate conditional expectations in the equilibrium conditions.
Conditional on this auxiliary object (and on a discrete \emph{localization index} that encodes the relevant piece of the state space or the binding pattern of constraints), the model is linear-Gaussian: the one-step transition is Gaussian with parameters given by the pathwise Jacobians, and the measurement equation is assumed linear with Gaussian noise.
Integrating out the auxiliary object yields a Gaussian-mixture transition kernel.
This structure enables a Rao--Blackwellized particle filter that marginalizes the high-dimensional continuous state with a Kalman filter and samples only the low-dimensional auxiliary object and discrete localization index.

We provide (i) a derivation of the conditional linear-Gaussian representation from the stacked equilibrium system via the implicit function theorem, (ii) sufficient conditions under which quadrature refinement converges to a well-defined mixture kernel and under which localization refinement yields a globally accurate piecewise-linear approximation, and (iii) stability-based bounds linking one-step kernel error to filtering and log-likelihood error.
Finally, we propose a benchmark New Keynesian model with an effective lower bound (ELB) on the nominal rate and outline a comparison to OccBin and to the piecewise-linear continuous (PLC) solution and conditionally optimal particle filter (COPF) of \citet{AruobaCubaBordaHigaFloresSchorfheideVillalvazo2021}.
\end{abstract}

\noindent\textbf{Keywords:} nonlinear DSGE estimation; particle filter; Rao--Blackwellization; extended path; stochastic extended path; dynamic perturbation; sigma points; occasionally binding constraints; effective lower bound.\\
\textbf{JEL:} C11, C15, C61, E32.

\section{Introduction}

Likelihood-based estimation of nonlinear DSGE models is difficult because the transition law for the latent state is both nonlinear and implicitly defined by equilibrium conditions with conditional expectations.
A ``global solution'' (projection, value function iteration, Chebyshev/Smolyak, endogenous grid, etc.) can be accurate but often becomes expensive as the dimension of the state grows, and embedding the global solution inside particle filtering and MCMC can be prohibitive.

This paper formalizes an alternative approximation strategy geared specifically toward likelihood evaluation:
instead of approximating the policy function globally, we approximate the \emph{local} dynamics in the neighborhood of \emph{paths} that are relevant under the data.
At each filtering step we solve for a finite-horizon equilibrium path under perfect foresight given a candidate sequence of future shocks, and then compute a first-order expansion around this path.
This is the spirit of the Extended Path (EP) method of \citet{FairTaylor1983} and its stochastic generalizations \citep{AdjemianJuillard2025}, and is closely related to Dynamic Perturbation \citep{MennuniRubioRamirezStepanchuk2025}, which computes first-order expansions along an entire simulated equilibrium path.

The novelty here is to place this pathwise approximation inside a \emph{filter}:
we construct a time-varying linear-Gaussian approximation \emph{conditional} on an auxiliary ``anticipated-shock'' object and (when needed) a discrete localization index that represents a regime/binding pattern or a partition cell in state space.
We then integrate out the auxiliary object by sigma points/quadrature, yielding a Gaussian-mixture kernel.
This makes the model a Gaussian-sum state-space system \citep{AlspachSorenson1972}, which can be handled with a Rao--Blackwellized particle filter (RBPF) \citep{DoucetFreitasMurphyRussell2000}.
With linear-Gaussian measurement equations, the RBPF marginalizes the large continuous state block by Kalman filtering conditional on the auxiliary variables.

A central question, raised by Ed's motivating idea, is:
\emph{when can such a method be ``better'' than either a global nonlinear solution or a bootstrap particle filter?}
The answer depends on an explicit bias--variance--cost tradeoff.
Pathwise linearization introduces approximation bias, but (i) it can be cheap relative to a global solution in large models because it relies on sparse Jacobians and linear solves, and (ii) Rao--Blackwellization can dramatically reduce Monte Carlo variance of likelihood estimates relative to a bootstrap particle filter in high-dimensional latent state spaces.

Finally, a second question is convergence:
\emph{as we add more approximation points, do we converge to something like a global solution?}
Sigma-point refinement alone converges to a \emph{mixture-of-linearizations} kernel, not necessarily to the true kernel.
Convergence to the true transition requires a refinement scheme that also shrinks the local-linearization and horizon/expectation-approximation errors.
We provide one such scheme based on a discrete localization index that refines a partition of the state space or the binding pattern of constraints.

\section{Nonlinear state-space model}

\subsection{Latent state and measurement}
Let $(x_t)_{t\ge 0}$ be an $\R^{n_x}$-valued latent state and $(y_t)_{t\ge 1}$ observed data.
We assume a linear measurement equation with Gaussian noise:
\begin{equation}
y_t = H x_t + \eta_t,\qquad \eta_t\sim \mathcal{N}(0,R),
\label{eq:meas}
\end{equation}
where $H$ and $R$ may depend on parameters $\theta$.

\subsection{Equilibrium transition and Gaussian shocks}
The state transition is induced by DSGE equilibrium conditions and structural innovations.
For exposition we write it as a Markov mapping
\begin{equation}
x_{t+1} = F(x_t,\varepsilon_{t+1};\theta),\qquad \varepsilon_{t+1}\sim \mathcal{N}(0,\Sigma),
\label{eq:true_transition}
\end{equation}
even though in DSGEs $F$ is usually not available in closed form.

Let $K_\theta(x,\cdot)$ denote the true Markov kernel of $x_{t+1}\mid x_t=x$ induced by \eqref{eq:true_transition}.

\section{From stacked equilibrium conditions to a conditional linear-Gaussian law}

This section derives the object the filter will use: a Gaussian approximation to $K_\theta(x_t,\cdot)$ conditional on an auxiliary object (anticipated shocks) and a discrete localization index.

\subsection{A finite-horizon ``stochastic perfect foresight'' system}

Fix a horizon $H\ge 2$.
Let
\[
\omega_t = (\varepsilon_{t+1}^a,\varepsilon_{t+2}^a,\ldots,\varepsilon_{t+H-1}^a)\in \R^{d},
\qquad d=(H-1)n_\varepsilon,
\]
collect a sequence of \emph{anticipated} shocks over the horizon.
We treat $\omega_t$ as an auxiliary variable drawn from the Gaussian law induced by the shock process:
\[
\omega_t \sim \mathcal{N}(0,\Omega),\qquad \Omega=\mathrm{diag}(\Sigma,\ldots,\Sigma).
\]
Conceptually, $\omega_t$ is a numerical device used to approximate conditional expectations, not an additional structural shock.

Let $X_{t:t+H-1}=(x_t,x_{t+1},\ldots,x_{t+H-1})$ and define a stacked residual mapping
\[
\mathcal{R}(X_{t:t+H-1};x_t,\omega_t,\theta)=0
\]
encoding equilibrium conditions from $t$ through $t+H-2$ plus a terminal condition at $t+H-1$ (e.g.\ $x_{t+H-1}\approx x^\star$).
In practice this is the system solved in EP/SEP and in perfect-foresight solvers \citep{FairTaylor1983,AdjemianJuillard2025}.
We denote by
\[
\bar X_{t:t+H-1}=\bar X(x_t,\omega_t;\theta)
\]
a solution of the stacked system.

\subsection{Local perturbations around the path}

Define deviations $\tilde x_s = x_s-\bar x_s$ for $s=t,\ldots,t+H-1$.
Linearizing $\mathcal{R}$ at $\bar X$ yields a sparse block system linking $\tilde x_{t:t+H-1}$ to \emph{unexpected} innovations.
To keep the derivation transparent, we focus on the one-step mapping from $(x_t,\varepsilon_{t+1})$ to $x_{t+1}$.
Let $\varepsilon_{t+1}$ be the \emph{unexpected} shock realization at $t+1$ (mean zero).
Then the linearized system implies
\begin{equation}
\tilde x_{t+1}
=
A_t(x_t,\omega_t)\,\tilde x_t
+
B_t(x_t,\omega_t)\,\varepsilon_{t+1},
\label{eq:lin_dev}
\end{equation}
where $A_t,B_t$ are obtained from the Jacobian of the stacked system (details in Appendix~\ref{app:ift}).

Equation \eqref{eq:lin_dev} yields a \emph{conditional} linear-Gaussian transition:
\begin{equation}
x_{t+1}\mid (x_t,\omega_t)
\;\approx\;
\mathcal{N}\!\left(m(x_t,\omega_t),\,Q(x_t,\omega_t)\right),
\label{eq:cond_gauss}
\end{equation}
with
\[
m(x_t,\omega_t)=\bar x_{t+1}+A_t(x_t,\omega_t)\,(x_t-\bar x_t),
\qquad
Q(x_t,\omega_t)=B_t(x_t,\omega_t)\,\Sigma\,B_t(x_t,\omega_t)^\top.
\]

\begin{remark}[Relation to dynamic perturbation and SEP]
Dynamic perturbation \citep{MennuniRubioRamirezStepanchuk2025} also computes first-order expansions along equilibrium paths.
SEP \citep{AdjemianJuillard2025} integrates over future shocks to better approximate conditional expectations than certainty equivalence.
The PPPF construction can be viewed as: (i) using a SEP-style draw/quadrature node $\omega_t$ to build the reference path, (ii) computing a dynamic-perturbation-style linearization around that path, and (iii) integrating out $\omega_t$ inside a filter.
\end{remark}

\section{Integrating out anticipated shocks}

\subsection{Gaussian-mixture transition kernel}

Integrating out $\omega_t$ yields an approximate transition kernel
\begin{equation}
\widehat K_\theta(x_t,\dd x_{t+1})
=
\int
\mathcal{N}\!\left(\dd x_{t+1};\,m(x_t,\omega),Q(x_t,\omega)\right)\,
\phi(\omega;0,\Omega)\,\dd\omega.
\label{eq:mixture_kernel}
\end{equation}
This is a Gaussian mixture (Gaussian-sum) transition \citep{AlspachSorenson1972}.
In practice we replace the integral by a $K$-node rule (sigma points, Gauss--Hermite, sparse grids, quasi-Monte Carlo, or Monte Carlo):
\begin{equation}
\widehat K_{\theta,K}(x_t,\dd x_{t+1})
=
\sum_{k=1}^K w_k\,
\mathcal{N}\!\left(\dd x_{t+1};\,m(x_t,\omega^{(k)}),Q(x_t,\omega^{(k)})\right),
\qquad \sum_{k=1}^K w_k=1.
\label{eq:quadrature_kernel}
\end{equation}

\subsection{Sigma points as deterministic ``anticipated shock'' draws}

For $\omega\sim \mathcal{N}(0,\Omega)\in\R^d$, the basic unscented transform uses $2d+1$ sigma points and associated weights that match mean and covariance and are exact for polynomials up to degree three (for symmetric choices) \citep{JulierUhlmann2004}.
Sigma points are attractive when $d$ is small (e.g.\ a small number of shocks and short horizon).

\section{Discrete localization and ``globalization'' of the approximation}

Sigma-point refinement controls the numerical integration error in \eqref{eq:mixture_kernel}, but does \emph{not} by itself eliminate the local-linearization and horizon/terminal-condition errors.
This section introduces a second approximation axis that can yield a genuine global refinement scheme.

\subsection{Localization index: partition cells or regime/binding patterns}

Let $s_t\in\{1,\ldots,J\}$ be a discrete localization index.
Two cases are especially relevant:

\begin{enumerate}[leftmargin=*, itemsep=2pt]
\item \textbf{State-space partition.}
Partition a compact set $\mathcal{X}\subset\R^{n_x}$ into cells $\{\mathcal{X}_j\}_{j=1}^J$ and choose anchors $x^{(j)}\in\mathcal{X}_j$.
For each anchor (and each quadrature node $\omega^{(k)}$) compute and store $(\bar x_{t+1}^{(j,k)},A^{(j,k)},B^{(j,k)})$.
At runtime, map each $x_t$ to its cell index $s_t=j$ and use the corresponding precomputed linearization parameters.
This turns online path solving into offline precomputation plus fast lookup/interpolation.

\item \textbf{Occasionally binding constraints (OBCs).}
Let $s_t$ encode the binding/slack pattern of a constraint over the horizon (a regime path).
OccBin \citep{GuerrieriIacoviello2015} selects a regime path by iterating on guesses until the constraint is satisfied.
In a filter, the regime can be treated as latent: particles can branch over regimes and the filter integrates over them.
Aruoba et al.\ \citep{AruobaCubaBordaHigaFloresSchorfheideVillalvazo2021} construct globally piecewise-linear and continuous decision rules and derive a conditionally optimal particle filter (COPF) that exploits this structure.
Our framework nests the idea of treating the binding pattern as a discrete index that makes the conditional system linear.
\end{enumerate}

\subsection{A simple approximation theorem for localization refinement}

To formalize ``convergence to global'' under partition refinement, we state a generic result.
For simplicity, consider a setting where $(m(x,\omega),Q(x,\omega))$ is smooth in $x$ for each $\omega$ (away from kinks).
Let $\widehat m_J(x,\omega)$ and $\widehat Q_J(x,\omega)$ be piecewise-linear approximations built from anchors with mesh size
\[
\Delta_J = \max_{j\le J} \sup_{x\in\mathcal{X}_j}\|x-x^{(j)}\|.
\]

\begin{assumption}[Smoothness in $x$]
For each $\omega$ in the support of the quadrature rule and all $x\in\mathcal{X}$, the maps $x\mapsto m(x,\omega)$ and $x\mapsto Q(x,\omega)$ are twice continuously differentiable with uniformly bounded second derivatives.
\end{assumption}

\begin{proposition}[Localization refinement implies uniform parameter convergence]
Under the smoothness assumption, there exists $C<\infty$ such that
\[
\sup_{x\in\mathcal{X}} \|m(x,\omega)-\widehat m_J(x,\omega)\|\le C \Delta_J^2,
\qquad
\sup_{x\in\mathcal{X}} \|Q(x,\omega)-\widehat Q_J(x,\omega)\|\le C \Delta_J^2,
\]
for each fixed $\omega$, provided $\widehat m_J,\widehat Q_J$ are constructed by first-order Taylor (or equivalent local linear) approximations on each cell.
Hence if $\Delta_J\to 0$ as $J\to\infty$, then the approximation error vanishes uniformly.
\end{proposition}

\begin{remark}
This is the standard ``mesh $\to 0$'' globalization logic: piecewise linearization converges for smooth maps.
For OBCs, the map is typically only piecewise smooth and the rate can degrade near the kink.
In that case one refines the partition more aggressively near the kink set, or treats the kink via an explicit discrete regime as in OccBin and PLC/COPF.
\end{remark}

\section{Filtering and likelihood evaluation}

\subsection{Gaussian-sum representation}

Under \eqref{eq:quadrature_kernel}, the approximate transition can be written by introducing a discrete index $k_t\in\{1,\ldots,K\}$ with $\Pr(k_t=k)=w_k$:
\[
x_{t+1}\mid (x_t,k_t=k) \sim \mathcal{N}(m_k(x_t),Q_k(x_t)).
\]
With localization $s_t$, the index becomes $(s_t,k_t)$.

\subsection{Rao--Blackwellized particle filter}

Because the measurement is linear-Gaussian and the transition is conditionally Gaussian given $(s_t,k_t)$, we can marginalize the continuous state by Kalman filtering conditional on each particle's discrete history.
A Rao--Blackwellized particle filter (RBPF) samples $(s_t,k_t)$ (or, in a continuous version, samples $\omega_t$) and updates the conditional Gaussian state distribution analytically \citep{DoucetFreitasMurphyRussell2000}.

\subsection{Algorithmic detail: mixture handling}

A key practical point is \emph{mixture handling}.
Given a mixture transition $\sum_k w_k \mathcal{N}(m_k,Q_k)$, averaging the means $m_k$ and covariances $Q_k$ to obtain a single Gaussian is generally incorrect; it destroys tail behavior and Jensen effects.
Instead, one must treat the mixture either by branching (expanding particles) or by sampling $k$ and accounting for the mixture density in the proposal.

In the RBPF setting, branching corresponds to spawning children for each $(s,k)$ and then resampling.
Sampling corresponds to drawing $(s,k)$ from their prior weights (or an auxiliary proposal) and updating particle weights accordingly.

\section{Approximation theory: kernel, filter, and likelihood}

\subsection{Quadrature refinement converges to the mixture kernel}

Let $\widehat K_\theta$ be the mixture kernel \eqref{eq:mixture_kernel} and $\widehat K_{\theta,K}$ its $K$-node quadrature approximation \eqref{eq:quadrature_kernel}.

\begin{assumption}[Quadrature consistency]
The quadrature rules $(\omega^{(k)},w_k)_{k=1}^K$ are consistent for the Gaussian measure $\mathcal{N}(0,\Omega)$:
\[
\sum_{k=1}^K w_k h(\omega^{(k)}) \to \int h(\omega)\,\phi(\omega;0,\Omega)\,\dd\omega
\quad \text{for all bounded continuous } h.
\]
\end{assumption}

\begin{proposition}[Weak convergence of $\widehat K_{\theta,K}\to \widehat K_\theta$]
Assume $(m(x,\omega),Q(x,\omega))$ is continuous in $\omega$ for each fixed $x$ and that dominated convergence applies to the Gaussian density integrand.
Then for each fixed $x$ and each bounded continuous test function $\psi$,
\[
\int \psi(x')\,\widehat K_{\theta,K}(x,\dd x') \to \int \psi(x')\,\widehat K_{\theta}(x,\dd x').
\]
\end{proposition}

\begin{remark}
This convergence holds to the mixture kernel $\widehat K_\theta$, not to the true DSGE kernel $K_\theta$.
Closing the remaining gap requires controlling the local-linearization error, horizon truncation/terminal-condition error, and expectation-approximation error in the underlying stochastic perfect-foresight step.
\end{remark}

\subsection{Kernel error decomposition}

Write the total kernel error as
\[
\|K_\theta(x,\cdot)-\widehat K_{\theta,K,J}(x,\cdot)\|
\;\le\;
\underbrace{\|K_\theta-\widetilde K_{\theta,H,p}\|}_{\text{horizon/expectation approximation}}
+
\underbrace{\|\widetilde K_{\theta,H,p}-\widehat K_{\theta}\|}_{\text{local linearization}}
+
\underbrace{\|\widehat K_{\theta}-\widehat K_{\theta,K}\|}_{\text{quadrature}}
+
\underbrace{\|\widehat K_{\theta,K}-\widehat K_{\theta,K,J}\|}_{\text{localization}},
\]
where $\widetilde K_{\theta,H,p}$ denotes the kernel induced by the \emph{exact} finite-horizon stochastic perfect-foresight solution with $p$ periods of expectation integration (no linearization), and $\widehat K_{\theta,K,J}$ denotes the practical PPPF kernel with $K$ quadrature nodes and localization mesh index $J$.

\subsection{From kernel error to filter and likelihood error}

Let $\delta=\sup_{x\in\mathcal{X}} \|K_\theta(x,\cdot)-\widehat K_{\theta,K,J}(x,\cdot)\|_{\mathrm{TV}}$ be the worst-case one-step kernel error.
Under standard mixing/forgetting conditions for hidden Markov models, errors in the transition kernel propagate to errors in filtering distributions and predictive likelihoods; see, e.g., \citet{DelMoral2004,Whiteley2013,KantasDoucetSinghMaciejowski2015}.

\begin{assumption}[Uniform mixing and bounded likelihood]
There exists $\epsilon\in(0,1)$ and a probability measure $\nu$ on $\mathcal{X}$ such that for all $x\in\mathcal{X}$ and measurable $A$,
\[
K_\theta(x,A)\ge \epsilon \nu(A),
\]
and the measurement density satisfies $0< g_{\min}\le g_\theta(y\mid x)\le g_{\max}<\infty$ for all $x\in\mathcal{X}$ and realized $y_t$.
\end{assumption}

\begin{theorem}[Filter perturbation bound (standard form)]
Under the mixing and bounded-likelihood assumption, there exist constants $C<\infty$ and $\rho\in(0,1)$ such that the filtering distributions satisfy
\[
\|\pi_{\theta,t}-\widehat \pi_{\theta,t}\|_{\mathrm{TV}}
\le
C \sum_{s=1}^{t} \rho^{t-s}\, \delta,
\qquad t\ge 1.
\]
\end{theorem}

\begin{corollary}[Log-likelihood perturbation bound (sketch)]
Under the same conditions, there exists $C_\ell<\infty$ such that
\[
\left| \log p_\theta(y_{1:T}) - \log \widehat p_\theta(y_{1:T}) \right|
\le
C_\ell\, T\, \delta
\]
up to geometric forgetting factors.
\end{corollary}

\section{When can PPPF be ``better'' than a global solution?}

Let $\ell_\theta$ denote the \emph{true} log-likelihood and let $\widehat \ell_\theta^{(A)}$ be a particle-based estimator using approximation strategy $A\in\{\mathrm{glob},\mathrm{pppf}\}$.
Under a compute budget $C$, the feasible particle count is roughly $N_A \approx C/c_A$, where $c_A$ is cost per particle per period.

A standard decomposition is
\[
\E\left[(\widehat \ell_\theta^{(A)}-\ell_\theta)^2\right]
=
\underbrace{\big(\mathrm{Bias}_A\big)^2}_{\text{approximation bias}}
+
\underbrace{\frac{\sigma_A^2}{N_A}}_{\text{Monte Carlo variance}}.
\]
Rao--Blackwellization reduces $\sigma_{\mathrm{pppf}}^2$ by integrating out (part of) the latent state analytically \citep{DoucetFreitasMurphyRussell2000}.
A sufficient ``dominance'' condition is therefore:
\[
(\mathrm{Bias}_{\mathrm{pppf}})^2 + \sigma_{\mathrm{pppf}}^2 \frac{c_{\mathrm{pppf}}}{C}
<
(\mathrm{Bias}_{\mathrm{glob}})^2 + \sigma_{\mathrm{glob}}^2 \frac{c_{\mathrm{glob}}}{C}.
\]
This inequality makes explicit why PPPF can lose in small models (high $c_{\mathrm{pppf}}$, weak RB gains) but win in large models where the global solution is expensive and RB gains are large.

\section{Benchmark: a New Keynesian model with an ELB}

This section specifies a concrete benchmark that directly connects to the OccBin and PLC/COPF literatures.

\subsection{Core equilibrium conditions}

We consider a canonical three-equation New Keynesian model with an effective lower bound (ELB) on the nominal rate.
Let $x_t$ denote the output gap, $\pi_t$ inflation, and $i_t$ the nominal interest rate (deviations from steady state, appropriately scaled).
A standard log-linearized NK system is
\begin{align}
x_t &= \E_t x_{t+1} - \sigma^{-1}\Big(i_t - \E_t \pi_{t+1} - r_t^n\Big), \label{eq:nk_is}\\
\pi_t &= \beta \E_t \pi_{t+1} + \kappa x_t, \label{eq:nk_pc}\\
i_t^\star &= \phi_\pi \pi_t + \phi_x x_t + \nu_t, \label{eq:taylor}\\
i_t &= \max\{ \underline i,\; i_t^\star\}. \label{eq:elb}
\end{align}
Here $r_t^n$ is a natural-rate (or preference) shock and $\nu_t$ a monetary policy shock; both follow AR(1) laws with Gaussian innovations.
Appendix~\ref{app:nk} provides a brief microfoundation and discusses common variations (habit, indexation, wage rigidities).

The occasionally binding ELB constraint \eqref{eq:elb} introduces a kink and makes the decision rules nonlinear and piecewise smooth.

\subsection{State-space representation and measurement}

A minimal latent state can be
\[
x_t^{\mathrm{state}} = (x_t,\pi_t,i_t,r_t^n,\nu_t)^\top,
\]
with observables
\[
y_t = (x_t^{\mathrm{obs}},\pi_t^{\mathrm{obs}},i_t^{\mathrm{obs}})^\top,
\qquad
y_t = H x_t^{\mathrm{state}} + \eta_t.
\]
Because the ELB affects the transition, likelihood evaluation requires a nonlinear filter unless the solution imposes a structure that can be exploited.

\subsection{Competing methods}

We propose comparing three likelihood-evaluation strategies:

\begin{enumerate}[leftmargin=*, itemsep=2pt]
\item \textbf{OccBin (piecewise linear perturbation) + bootstrap PF / inversion filter.}
OccBin \citep{GuerrieriIacoviello2015} computes a piecewise linear solution by iterating on the binding pattern of \eqref{eq:elb} over a finite horizon.
Likelihood evaluation can be done by a bootstrap particle filter, and several papers use inversion-filter variants when the observables identify shocks (or approximately identify them).

\item \textbf{Piecewise-linear continuous (PLC) decision rules + COPF.}
\citet{AruobaCubaBordaHigaFloresSchorfheideVillalvazo2021} construct globally piecewise-linear and continuous decision rules for OBC models and derive a \emph{conditionally optimal} particle filter (COPF) that exploits the piecewise-linear structure to reduce likelihood variance.

\item \textbf{PPPF with anticipated-shock quadrature and regime localization.}
In the PPPF approach, we treat the ELB binding pattern over the EP horizon as part of the localization index $s_t$ and integrate over short-horizon anticipated shocks $\omega_t$ by sigma points/quadrature.
Conditional on $(s_t,\omega_t)$ the system is linear, enabling Rao--Blackwellization.
\end{enumerate}

\subsection{Proposed evaluation metrics}

Because particle filters yield stochastic likelihood approximations, the key objects for comparison are:
\begin{itemize}[leftmargin=*, itemsep=2pt]
\item runtime per likelihood evaluation,
\item Monte Carlo variance of the log-likelihood estimate (across repeated PF runs on fixed data),
\item effective sample size (ESS) profiles and frequency of resampling,
\item impulse responses: nonlinear ``true'' (high-accuracy) vs OccBin/PLC vs PPPF.
\end{itemize}
This benchmark is designed to highlight cases where PPPF can plausibly outperform a naive bootstrap PF: near the ELB kink, the conditional Gaussian structure and regime-aware localization can provide substantially better proposals and variance reduction.

\section{Conclusion}

PPPF replaces a global nonlinear approximation of DSGE policy functions with many local linearizations around stochastic perfect-foresight paths, integrating over an anticipated-shock object by sigma points/quadrature and optionally over regimes via a discrete localization index.
With Gaussian shocks and linear measurement equations, this yields a Gaussian-mixture state-space model amenable to Rao--Blackwellized particle filtering.
We provided derivations, convergence statements for quadrature and localization refinement, and stability-based links from kernel error to likelihood error.
A New Keynesian ELB benchmark offers a natural comparison to the OccBin and PLC/COPF literatures.

\appendix

\section{Deriving the conditional linear system via the implicit function theorem}
\label{app:ift}

This appendix provides a derivation of \eqref{eq:lin_dev} from the stacked residual system.
Let $Y$ denote the stacked unknowns over the horizon excluding the pinned initial state:
\[
Y = (x_{t+1},x_{t+2},\ldots,x_{t+H-1})\in \R^{(H-1)n_x}.
\]
Write the stacked equilibrium conditions as
\[
\mathcal{F}(Y;x_t,\omega_t,\theta)=0.
\]
Suppose $\bar Y$ solves the system for given $(x_t,\omega_t)$.
Assume the Jacobian with respect to $Y$ is nonsingular:
\[
D_Y \mathcal{F}(\bar Y;x_t,\omega_t,\theta)\;\;\text{is invertible}.
\]
Then by the implicit function theorem there exists a local mapping $Y=G(x_t,\omega_t)$ with derivatives
\[
D_{x}G = -\big(D_Y\mathcal{F}\big)^{-1} D_x \mathcal{F},
\qquad
D_{\omega}G = -\big(D_Y\mathcal{F}\big)^{-1} D_\omega \mathcal{F}.
\]
To incorporate an \emph{unexpected} shock $\varepsilon_{t+1}$, one augments the residual function to include the contemporaneous shock in the appropriate period-$t+1$ equations, and differentiates with respect to $\varepsilon_{t+1}$.
Because $D_Y\mathcal{F}$ is sparse block-banded (typically block tridiagonal for one-step-ahead expectations), solving for these derivatives amounts to solving a sparse linear system, which is computationally attractive in large DSGEs.

Extracting the first block row of $D_x G$ gives the mapping from $x_t$ perturbations to $x_{t+1}$ perturbations, yielding $A_t$.
Similarly, the first block row of $D_{\varepsilon} G$ gives $B_t$.
Plugging into $x_{t+1}=\bar x_{t+1}+\tilde x_{t+1}$ gives \eqref{eq:lin_dev}.

\section{Unscented-transform sigma points}
\label{app:ut}

For $\omega\sim\mathcal{N}(0,\Omega)\in\R^d$, the basic sigma-point set uses $2d+1$ points:
\[
\omega^{(0)}=0,\qquad
\omega^{(i)}=+\left[\sqrt{(d+\lambda)\Omega}\right]_{\cdot i},\qquad
\omega^{(i+d)}=-\left[\sqrt{(d+\lambda)\Omega}\right]_{\cdot i},
\quad i=1,\ldots,d,
\]
with weights
\[
w_0^m=\frac{\lambda}{d+\lambda},\qquad w_i^m=\frac{1}{2(d+\lambda)}\;\;(i\ge 1),
\]
and covariance weights $w_0^c=w_0^m + (1-\alpha^2+\beta)$, $w_i^c=w_i^m$.
Standard choices are $\alpha\in(0,1]$, $\beta=2$ (Gaussian prior), $\kappa=0$ \citep{JulierUhlmann2004}.

\section{Microfoundation and variants of the NK--ELB benchmark}
\label{app:nk}

Equations \eqref{eq:nk_is}--\eqref{eq:nk_pc} arise from household Euler equations and Calvo price setting under standard assumptions; see \citet{Woodford2003} for a textbook derivation.
The ELB \eqref{eq:elb} is an occasionally binding constraint.
In empirical work one often uses a shadow-rate specification or introduces additional wedges, but the essential computational feature is the kink introduced by the max operator.

\begin{thebibliography}{99}

\bibitem[Adjemian and Juillard(2025)]{AdjemianJuillard2025}
Adjemian, S. and M. Juillard (2025).
\newblock Stochastic extended path.
\newblock \emph{Dynare Working Papers} No.\ 84, revised September 2025.

\bibitem[Alspach and Sorenson(1972)]{AlspachSorenson1972}
Alspach, D. L. and H. W. Sorenson (1972).
\newblock Nonlinear Bayesian estimation using Gaussian sum approximations.
\newblock \emph{IEEE Transactions on Automatic Control} 17(4), 439--448.

\bibitem[Aruoba et~al.(2021)]{AruobaCubaBordaHigaFloresSchorfheideVillalvazo2021}
Aruoba, S. B., P. Cuba-Borda, K. Higa-Flores, F. Schorfheide, and S. Villalvazo (2021).
\newblock Piecewise-linear approximations and filtering for DSGE models with occasionally binding constraints.
\newblock \emph{Review of Economic Dynamics} 41, 96--120.
\newblock DOI: 10.1016/j.red.2020.12.003.

\bibitem[Del~Moral(2004)]{DelMoral2004}
Del~Moral, P. (2004).
\newblock \emph{Feynman--Kac Formulae: Genealogical and Interacting Particle Systems with Applications}.
\newblock Springer.

\bibitem[Doucet et~al.(2000)]{DoucetFreitasMurphyRussell2000}
Doucet, A., N. de~Freitas, K. Murphy, and S. Russell (2000).
\newblock Rao--Blackwellised particle filtering for dynamic Bayesian networks.
\newblock In \emph{Proceedings of the Sixteenth Conference on Uncertainty in Artificial Intelligence (UAI)}.

\bibitem[Fair and Taylor(1983)]{FairTaylor1983}
Fair, R. C. and J. B. Taylor (1983).
\newblock Solution and maximum likelihood estimation of dynamic nonlinear rational expectations models.
\newblock \emph{Econometrica} 51(4), 1169--1185.

\bibitem[Guerrieri and Iacoviello(2015)]{GuerrieriIacoviello2015}
Guerrieri, L. and M. Iacoviello (2015).
\newblock OccBin: A toolkit for solving dynamic models with occasionally binding constraints easily.
\newblock \emph{Journal of Monetary Economics} 70, 22--38.
\newblock DOI: 10.1016/j.jmoneco.2014.08.005.

\bibitem[Gust et~al.(2017)]{GustHerbstLopezSalidoSmith2017}
Gust, C., E. Herbst, D. L{\'o}pez-Salido, and M. E. Smith (2017).
\newblock The empirical implications of the interest-rate lower bound.
\newblock \emph{American Economic Review} 107(7), 1971--2006.
\newblock DOI: 10.1257/aer.20121437.

\bibitem[Herbst and Schorfheide(2019)]{HerbstSchorfheide2019}
Herbst, E. and F. Schorfheide (2019).
\newblock Tempered particle filtering.
\newblock \emph{Journal of Econometrics} 210(1), 26--44.

\bibitem[Julier and Uhlmann(2004)]{JulierUhlmann2004}
Julier, S. J. and J. K. Uhlmann (2004).
\newblock Unscented filtering and nonlinear estimation.
\newblock \emph{Proceedings of the IEEE} 92(3), 401--422.

\bibitem[Kantas et~al.(2015)]{KantasDoucetSinghMaciejowski2015}
Kantas, N., A. Doucet, S. Singh, and J. Maciejowski (2015).
\newblock On particle methods for parameter estimation in state-space models.
\newblock \emph{Statistical Science} 30(3), 328--351.

\bibitem[Mennuni et~al.(2025)]{MennuniRubioRamirezStepanchuk2025}
Mennuni, A., J. F. Rubio-Ram{\'\i}rez, and S. Stepanchuk (2025).
\newblock Dynamic perturbation.
\newblock \emph{The Review of Economic Studies} 92(2), 1157--1192.
\newblock DOI: 10.1093/restud/rdae037.

\bibitem[Whiteley(2013)]{Whiteley2013}
Whiteley, N. (2013).
\newblock Stability properties of some particle filters.
\newblock \emph{The Annals of Applied Probability} 23(6), 2500--2537.

\bibitem[Woodford(2003)]{Woodford2003}
Woodford, M. (2003).
\newblock \emph{Interest and Prices: Foundations of a Theory of Monetary Policy}.
\newblock Princeton University Press.

\end{thebibliography}

\end{document}
